\documentclass[10pt, aspectratio=169]{beamer}
% Preámbulo modular para presentaciones Beamer (Modelación Estadística)
% Diseño y paleta institucional del Tecnológico de Monterrey

\documentclass[aspectratio=169,xcolor={dvipsnames,table}]{beamer}

\usepackage[utf8]{inputenc}
\usepackage[spanish,mexico]{babel}
\usepackage{amsmath,amssymb,amsfonts,mathtools}
\usepackage{graphicx}
\usepackage{listings}
\usepackage{booktabs}
\usepackage{tikz}
\usetikzlibrary{matrix,arrows,decorations.pathmorphing,shapes.geometric,positioning}

% Paleta Institucional Tecnológico de Monterrey
\definecolor{TecAzul}{HTML}{0039A6}       % Azul TEC: color primario institucional
\definecolor{TecAzulOscuro}{HTML}{1A2E51} % Azul marino oscuro
\definecolor{TecRojo}{HTML}{EC2661}       % Rojo de acento (paleta miTec)
\definecolor{TecGris}{HTML}{646464}       % Gris neutro para comentarios y subtítulos
\definecolor{TecFondoBloque}{HTML}{F4F6F9}% Fondo suave para bloques

% Configuración de Tema Beamer
\usetheme{Boadilla}
\usecolortheme[named=TecAzulOscuro]{structure}
\setbeamercolor{alerted text}{fg=TecRojo}
\setbeamercolor{palette primary}{bg=TecAzulOscuro,fg=white}
\setbeamercolor{palette secondary}{bg=TecAzul,fg=white}
\setbeamercolor{palette tertiary}{bg=TecAzulOscuro!80!black,fg=white}
\setbeamercolor{title}{fg=TecAzulOscuro,bg=TecFondoBloque}
\setbeamercolor{frametitle}{fg=TecAzulOscuro,bg=TecFondoBloque}

% Bloques personalizados elegantes
\setbeamercolor{block title}{bg=TecAzulOscuro,fg=white}
\setbeamercolor{block body}{bg=TecFondoBloque,fg=black}
\setbeamercolor{block title alerted}{bg=TecRojo,fg=white}
\setbeamercolor{block body alerted}{bg=TecRojo!8,fg=black}
\setbeamercolor{block title example}{bg=TecAzul,fg=white}
\setbeamercolor{block body example}{bg=TecAzul!8,fg=black}

% Redondear bordes y añadir sombra sutil
\setbeamertemplate{blocks}[rounded][shadow=true]
\setbeamertemplate{navigation symbols}{}

% Pie de página limpio con número de diapositiva
\setbeamertemplate{footline}{
  \leavevmode%
  \hbox{%
  \begin{beamercolorbox}[wd=.7\paperwidth,ht=2.25ex,dp=1ex,left]{author in head/foot}%
    \hspace*{2ex}\insertshortauthor~~(\insertshortinstitute) \hfill \insertshorttitle
  \end{beamercolorbox}%
  \begin{beamercolorbox}[wd=.3\paperwidth,ht=2.25ex,dp=1ex,right]{date in head/foot}%
    \insertshortdate{}\hspace*{2em}
    \insertframenumber{} / \inserttotalframenumber\hspace*{2ex}
  \end{beamercolorbox}}%
  \vskip0pt%
}

% Configuración de Listings de Python para visualización pedagógica de código
\lstset{
    language=Python,
    basicstyle=\ttfamily\footnotesize,
    keywordstyle=\color{TecAzulOscuro}\bfseries,
    stringstyle=\color{TecRojo},
    commentstyle=\color{TecGris}\itshape,
    showstringspaces=false,
    breaklines=true,
    frame=lines,
    rulecolor=\color{TecAzulOscuro!30},
    backgroundcolor=\color{TecFondoBloque},
    aboveskip=0.5em,
    belowskip=0.5em,
    literate={á}{{\'a}}1 {é}{{\'e}}1 {í}{{\'i}}1 {ó}{{\'o}}1 {ú}{{\'u}}1
             {Á}{{\'A}}1 {É}{{\'E}}1 {Í}{{\'I}}1 {Ó}{{\'O}}1 {Ú}{{\'U}}1
             {ñ}{{\~n}}1 {Ñ}{{\~N}}1 {¿}{{?`}}1 {¡}{{!`}}1
}

% Metadatos institucionales por defecto
\title[Modelación Estadística]{Modelación Estadística para Ciencia de Datos e Ingeniería}
\author[J. Castillo Colmenares]{Juliho David Castillo Colmenares}
\institute[Tec de Monterrey]{Tecnológico de Monterrey \\ Escuela de Ingeniería y Ciencias}
\date{\today}

% Comandos y macros estadísticas compartidas para las presentaciones Beamer (Metropolis)

% Conjuntos numéricos básicos
\newcommand{\R}{\mathbb{R}}
\newcommand{\N}{\mathbb{N}}
\newcommand{\Q}{\mathbb{Q}}
\newcommand{\Z}{\mathbb{Z}}
\newcommand{\set}[1]{\left\{ #1 \right\}}
\newcommand{\abs}[1]{\left|#1\right|}
\newcommand{\norm}[1]{\left\| #1 \right\|}

% Comandos estadísticos y probabilísticos del curso
\newcommand{\Var}{\operatorname{Var}}
\newcommand{\cov}{\operatorname{Cov}}
\newcommand{\corr}{\rho}
\newcommand{\comb}[2]{\begin{pmatrix} #1 \\ #2 \end{pmatrix}}
\newcommand{\s}{\sigma}
\newcommand{\particion}{\mathfrak{P}}
\newcommand{\card}[1]{\#\left( #1 \right)}
\newcommand{\seq}[3]{\set{#1}_{#2}^{#3}}
\newcommand{\E}{\mathbb{E}}
\newcommand{\Prob}{\mathbb{P}}

% Entornos pedagógicos y cajas en español académico natural
\newenvironment{discusion}{%
  \begin{alertblock}{Pregunta de Discusión}%
}{%
  \end{alertblock}%
}

\newenvironment{reflexion}{%
  \begin{alertblock}{Reflexión para el Aula}%
}{%
  \end{alertblock}%
}

\newenvironment{paradoja}{%
  \begin{alertblock}{Paradoja de Exploración}%
}{%
  \end{alertblock}%
}

\newenvironment{motivacion}{%
  \begin{exampleblock}{Motivación: Ciencia de Datos en la Práctica}%
}{%
  \end{exampleblock}%
}

\newenvironment{retoclase}{%
  \begin{block}{Reto Rápido en Equipo}%
}{%
  \end{block}%
}


\title[Distribución Normal]{Sección 03.09: Distribución Normal}
\subtitle{Aproximación Continua de Variables Discretas y Teorema del Límite Central}
\author[J. Castillo Colmenares]{Juliho Castillo Colmenares}
\institute{Tecnológico de Monterrey}
\date{\vspace{-1.2cm}}

\begin{document}

% Slide 1: Portada
\begin{frame}[plain]
  \titlepage
\end{frame}

% Slide 2: Hoja de Ruta
\begin{frame}{Hoja de Ruta --- Capítulo 03: Variables Aleatorias Discretas}
  \footnotesize
  ¿Dónde nos encontramos en el desarrollo modular de la teoría? \pause
  \begin{itemize}\itemsep=0.08em
    \item \textbf{03.01} PMF y Soporte de Variables Aleatorias Discretas \pause
    \item \textbf{03.02} Función de Distribución Acumulada Discreta (CDF) \pause
    \item \textbf{03.03} Esperanza Matemática y Varianza \pause
    \item \textbf{03.04} Distribución Binomial y Bernoulli \pause
    \item \textbf{03.05} Distribuciones Geométrica y Binomial Negativa \pause
    \item \textbf{03.06} Distribución Hipergeométrica \pause
    \item \textbf{03.07} Distribución de Poisson \pause
    \item \textbf{03.08} Distribución Multinomial \pause
    \item \color{TecRojo}\textbf{03.09 Distribución Normal \emph{(Hoy)}}\color{black}
  \end{itemize}
  \vspace{-0.1cm}
  \begin{block}{Objetivo de la Sesión}
    Introducir la distribución Normal como modelo continuo fundamental, dominar la estandarización al puntaje $Z$ y la regla 68-95-99.7, explotar el Teorema de De Moivre-Laplace para aproximar la Binomial por la Normal con corrección de continuidad, y aplicar el Teorema Central del Límite a sumas de variables discretas.
  \end{block}
\end{frame}

% Slide 3: Motivación
\begin{frame}{De Discreto a Continuo: El Teorema del Límite como Puente}
  \footnotesize
  \begin{motivacion}[¿Por qué la distribución Normal?]
    Hasta ahora hemos estudiado distribuciones para variables aleatorias \emph{discretas} que toman valores enteros (Binomial, Poisson, Multinomial). Pero muchos fenómenos en ingeniería y ciencia varían continuamente: errores de medición, fluctuaciones de precios, ruidos eléctricos, tiempos entre eventos. La distribución Normal, con su característica forma de campana, es el modelo universal para estos fenómenos y el límite asintótico de casi todas las distribuciones discretas.
  \end{motivacion}

  \vspace{0.15cm}
  \pause
  \begin{itemize}\itemsep=0.1cm
    \item \textbf{Carácter universal:} Aparece en el TLC como límite de sumas de variables i.i.d. con media y varianza finitas. \pause
    \item \textbf{Modelo continuo:} Función de densidad $f(x)$ en lugar de masa de probabilidad $P(X=k)$; soporte en todo $\mathbb{R}$. \pause
    \item \textbf{Puente discreto-continuo:} Permite aproximar la Binomial y la Poisson cuando $n$ o $\lambda$ son grandes.
  \end{itemize}
\end{frame}

% Slide 4: Definición Formal
\begin{frame}{Definición Formal de la Distribución Normal}
  \footnotesize
  Una variable aleatoria continua $X$ sigue una distribución Normal con parámetros $\mu \in \mathbb{R}$ y $\sigma > 0$, denotada $X \sim N(\mu, \sigma^{2})$, si su función de densidad de probabilidad (PDF) es: \pause

  \vspace{0.1cm}
  \begin{block}{Función de Densidad de Probabilidad (PDF)}
    $$f(x) = \dfrac{1}{\sigma\sqrt{2\pi}} \exp\left(-\dfrac{(x - \mu)^{2}}{2\sigma^{2}}\right), \quad x \in \mathbb{R}$$
  \end{block}

  \vspace{0.1cm}
  \pause
  \begin{alertblock}{Interpretación Geométrica y Parametrización en SciPy}
    \begin{itemize}\itemsep=0.06cm
      \item \textbf{Campana centrada en $\mu$:} La densidad es máxima en $x = \mu$ con valor $1/(\sigma\sqrt{2\pi})$.
      \item \textbf{Dispersión controlada por $\sigma$:} Mayor $\sigma$ implica campana más ancha y baja.
      \item \textbf{En SciPy:} \texttt{norm.pdf(x, loc=mu, scale=sigma)} y \texttt{norm.cdf(x, loc=mu, scale=sigma)}.
    \end{itemize}
  \end{alertblock}
\end{frame}

% Slide 5: Estandarización
\begin{frame}{Estandarización: Puntaje $Z \sim N(0, 1)$}
  \footnotesize
  Cualquier variable Normal $X \sim N(\mu, \sigma^{2})$ puede transformarse en una Normal estándar $Z \sim N(0, 1)$ mediante: \pause

  \vspace{0.1cm}
  \begin{block}{Transformación de Estandarización}
    $$Z = \dfrac{X - \mu}{\sigma} \implies Z \sim N(0, 1)$$
    La inversa: $X = \mu + \sigma Z$. La estandarización preserva las probabilidades: $P(a \le X \le b) = P\left(\dfrac{a - \mu}{\sigma} \le Z \le \dfrac{b - \mu}{\sigma}\right)$.
  \end{block}

  \vspace{0.1cm}
  \pause
  \begin{alertblock}{CDF Normal Estándar y su Simetría}
    \begin{itemize}\itemsep=0.05cm
      \item $\Phi(z) = P(Z \le z) = \int_{-\infty}^{z} \dfrac{1}{\sqrt{2\pi}} e^{-u^{2}/2}\,du$.
      \item $\Phi(-z) = 1 - \Phi(z)$ (simetría de la campana).
      \item Valores tabulados: $\Phi(0) = 0.5$, $\Phi(1) \approx 0.8413$, $\Phi(2) \approx 0.9772$.
    \end{itemize}
  \end{alertblock}
\end{frame}

% Slide 6: Regla 68-95-99.7
\begin{frame}{Regla Empírica 68-95-99.7 y Propiedades de Simetría}
  \footnotesize
  La distribución Normal tiene una regla mnemotécnica fundamental para interpretar rápidamente probabilidades: \pause

  \vspace{0.1cm}
  \begin{block}{Regla Empírica (Empirical Rule)}
    \begin{itemize}\itemsep=0.06cm
      \item $P(\mu - \sigma \le X \le \mu + \sigma) = P(|Z| \le 1) = 2\Phi(1) - 1 \approx 0.6827$ (68\%)
      \item $P(\mu - 2\sigma \le X \le \mu + 2\sigma) = P(|Z| \le 2) = 2\Phi(2) - 1 \approx 0.9545$ (95\%)
      \item $P(\mu - 3\sigma \le X \le \mu + 3\sigma) = P(|Z| \le 3) = 2\Phi(3) - 1 \approx 0.9973$ (99.7\%)
    \end{itemize}
  \end{block}

  \vspace{0.1cm}
  \pause
  \begin{alertblock}{Aplicación en Control de Calidad}
    Las cartas de control $\bar{x}$ usan límites $\mu \pm 3\sigma$ que capturan $99.73\%$ de la variación bajo operación normal. Puntos fuera de estos límites señalan ``causas especiales'' de variabilidad que requieren investigación.
  \end{alertblock}
\end{frame}

% Slide 7: Teorema de De Moivre-Laplace
\begin{frame}{Teorema de De Moivre-Laplace: Binomial $\to$ Normal}
  \footnotesize
  Cuando $n$ es grande, la Binomial $X \sim \text{Bin}(n, p)$ se aproxima por la Normal con los mismos parámetros: \pause

  \vspace{0.1cm}
  \begin{block}{Límite Asintótico de la Binomial}
    $$\dfrac{X - np}{\sqrt{np(1-p)}} \xrightarrow{d} N(0, 1) \quad \text{cuando } n \to \infty$$
    La aproximación es excelente bajo la regla operativa $np \ge 5$ y $n(1-p) \ge 5$. Para mejorar la precisión, se aplica la \emph{corrección de continuidad de Yates}.
  \end{block}

  \vspace{0.1cm}
  \pause
  \begin{alertblock}{Corrección de Continuidad de Yates}
    Para aproximar $P(X = k)$ o $P(X \ge k)$ se usa $P(k - 0.5 \le Y \le k + 0.5)$ con $Y \sim N(np, np(1-p))$. Esto reduce el error sistemático de tratar una variable discreta como continua.
  \end{alertblock}
\end{frame}

% Slide 8: Aproximación Poisson-Normal
\begin{frame}{Aproximación Poisson $\to$ Normal}
  \footnotesize
  La distribución de Poisson $X \sim \text{Pois}(\lambda)$ también converge a la Normal cuando $\lambda$ crece: \pause

  \vspace{0.1cm}
  \begin{block}{Límite Asintótico de Poisson}
    $$\dfrac{X - \lambda}{\sqrt{\lambda}} \xrightarrow{d} N(0, 1) \quad \text{cuando } \lambda \to \infty$$
    La regla operativa es $\lambda \ge 30$ para que la asimetría de la Poisson sea despreciable. La aproximación es fundamental en ingeniería de tráfico y teoría de colas.
  \end{block}

  \vspace{0.1cm}
  \pause
  \begin{alertblock}{Aplicación en Seguros y Salud Pública}
    El número diario de reclamaciones a una aseguradora ($\lambda = 100$) o el número semanal de casos de una enfermedad ($\lambda = 50$) se modelan eficientemente con $N(\lambda, \lambda)$, simplificando el cálculo de probabilidades de colas y percentiles operativos.
  \end{alertblock}
\end{frame}

% Slide 9: Teorema Central del Límite
\begin{frame}{Teorema Central del Límite: Universalidad de la Normal}
  \footnotesize
  El Teorema Central del Límite (TLC) establece que la suma de variables aleatorias i.i.d. con media y varianza finitas converge a la Normal: \pause

  \vspace{0.1cm}
  \begin{block}{Teorema Central del Límite (TLC)}
    Sea $X_{1}, X_{2}, \dots, X_{n}$ una muestra i.i.d. con media $\mu$ y varianza $\sigma^{2} < \infty$. Entonces:
    $$Z_{n} = \dfrac{\sum_{i=1}^{n} X_{i} - n\mu}{\sigma\sqrt{n}} \xrightarrow{d} N(0, 1) \quad \text{cuando } n \to \infty$$
  \end{block}

  \vspace{0.1cm}
  \pause
  \begin{alertblock}{Aplicación para $n = 30$}
    La calidad de la aproximación mejora con $n$. Para $n = 30$, la aproximación es típicamente excelente para distribuciones simétricas; para $n \ge 50$ funciona bien incluso con distribuciones sesgadas como la exponencial.
  \end{alertblock}
\end{frame}

% Slide 10: Aplicaciones
\begin{frame}{Aplicaciones en Ingeniería, Ciencias y Ciencia de Datos}
  \footnotesize
  La distribución Normal es omnipresente en aplicaciones cuantitativas: \pause

  \vspace{0.15cm}
  \begin{itemize}\itemsep=0.15cm
    \item \textbf{Control de Calidad Industrial:} Monitoreo de procesos de manufactura con cartas de control $\bar{x}$ y $R$ basadas en $N(\mu, \sigma^{2})$. \pause
    \item \textbf{Inferencia Estadística:} Pruebas $Z$, intervalos de confianza, y base teórica del análisis de varianza (ANOVA) y regresión lineal. \pause
    \item \textbf{Física e Instrumentación:} Errores de medición, ruido gaussiano en sistemas de comunicación, y propagación de incertidumbres. \pause
    \item \textbf{Finanzas y Análisis de Riesgo:} Modelo de Black-Scholes para precios de opciones y medida de valor en riesgo (VaR).
  \end{itemize}
\end{frame}

% Slide 12A-1: Lab Python (Bloque 1)
\begin{frame}[fragile]{Laboratorio en Python: PDF Normal y Estandarización Z}
  \scriptsize
  Verificación computacional en SciPy (\texttt{03.09\_normal\_approximation.py}) de la PDF, normalización y estandarización al puntaje $Z$:
  \vspace{0.05cm}
  {\lstset{basicstyle=\fontsize{5pt}{6pt}\selectfont\ttfamily}
  \lstinputlisting[firstline=15, lastline=37]{../../code/03_variables_aleatorias_discretas/03.09_normal_approximation.py}}
\end{frame}

% Slide 12A-2: Lab Python (Bloque 2)
\begin{frame}[fragile]{Laboratorio en Python: Aproximación Binomial-Normal}
  \scriptsize
  De Moivre-Laplace: comparación de PMF Binomial con PDF Normal con corrección de continuidad de Yates:
  \vspace{0.02cm}
  {\lstset{basicstyle=\fontsize{4.5pt}{5.5pt}\selectfont\ttfamily}
  \lstinputlisting[firstline=40, lastline=64]{../../code/03_variables_aleatorias_discretas/03.09_normal_approximation.py}}
\end{frame}

% Slide 12B: Lab Python (Bloque 3)
\begin{frame}[fragile]{Laboratorio en Python: Poisson-Normal y TLC}
  \scriptsize
  Aproximación Poisson $\to$ Normal y verificación Monte Carlo del TLC para suma de 50 exponenciales:
  \vspace{0.02cm}
  {\lstset{basicstyle=\fontsize{4.5pt}{5.5pt}\selectfont\ttfamily}
  \lstinputlisting[firstline=73, lastline=100]{../../code/03_variables_aleatorias_discretas/03.09_normal_approximation.py}}
\end{frame}

% Slide 12C: Salida Terminal
\begin{frame}[fragile]{Laboratorio en Python: Salida en Terminal}
  \scriptsize
  Salida estándar generada al ejecutar el script del laboratorio en Python:
  \vspace{0.02cm}
  \begin{block}{Salida Estándar en Consola}
    \fontsize{4pt}{5pt}\selectfont
    \begin{verbatim}
=== Block 1: Normal PDF & Z-Score Standardization ===
X ~ N(2.50, 0.05) | PDF integral: 1.000000
P(2.42 <= X <= 2.58) exact = 0.890401 | 2*Phi(1.60)-1 = 0.890401
Empirical rule: P(|Z|<=1)=0.6827 | P(|Z|<=2)=0.9545 | P(|Z|<=3)=0.9973

=== Block 2: Binomial-to-Normal Approximation ===
Bin(200, 0.25): mu=50.0, sigma=6.1237
  P(X>=60) exact Binomial:  0.062472
  P(X>=60) with Yates:     0.060410 (Z=1.5513)
  P(X>=60) without Yates:  0.051235 (Z=1.6330)
  Error with Yates:    0.002063
  Error without Yates: 0.011237
Convergence: n=50: exact=0.556138 | normal=0.556231

=== Block 3: Poisson-to-Normal & CLT ===
Pois(100) -> N(100, 100): P(85<=X<=115) exact=0.879276
  Normal approx: 0.878858 | Error: 0.000417
Sum of 100 Pois(2) -> S_100 ~ Pois(200):
  P(S_100<=220) exact=0.924699 | normal=0.926411
Sum of 50 Exp(1) ~ Gamma(50, 1):
  P(45<=S_50<=55) exact=0.520993 | normal=0.520500
Monte Carlo (N=250,000): mean=49.9981 | std=7.0764
  P(45<=S<=55) empirical=0.519140
    \end{verbatim}
  \end{block}
\end{frame}

% Slide 13A: Ejercicio Nivel 1 (Enunciado)

% Slide 13B: Ejercicio Nivel 1 (Resolución)

% Slide 14A: Ejercicio Nivel 2 (Enunciado)

% Slide 14B: Ejercicio Nivel 2 (Resolución)

% Slide 15A: Ejercicio Nivel 3 (Enunciado)

% Slide 17: Cierre
\begin{frame}{Cierre de Sección y Síntesis sobre la Distribución Normal}
  \begin{reflexion}[La Universalidad de la Campana de Gauss]
    La distribución Normal emerge como el modelo continuo fundamental que cierra el puente entre lo discreto y lo continuo. Desde el Teorema de De Moivre-Laplace hasta el TLC, la campana de Gauss aparece como el destino asintótico universal de las distribuciones de conteo cuando los parámetros de escala crecen suficientemente.
  \end{reflexion}

  \vspace{0.2cm}
  \pause
  \begin{block}{Lecciones Clave de la Distribución Normal}
    \begin{itemize}\itemsep=0.08cm
      \item \textbf{Carácter universal:} La Normal es el límite de la Binomial, Poisson y suma de variables i.i.d. con varianza finita.
      \item \textbf{Estandarización $Z$:} Cualquier cálculo se reduce a la tabla $\Phi(\cdot)$ mediante $Z = (X - \mu)/\sigma$.
      \item \textbf{Corrección de Yates:} Esencial para reducir el error sistemático al aproximar distribuciones discretas por la Normal.
    \end{itemize}
  \end{block}
\end{frame}

% Slide 18: Perspectiva Modular
\begin{frame}{Perspectiva Modular --- El Próximo Reto}
  \scriptsize
  \begin{retoclase}[Para la próxima sesión: Distribuciones Discretas en Ciencia de Datos]
    Hemos dominado las distribuciones discretas fundamentales (Binomial, Geométrica, Hipergeométrica, Poisson, Multinomial) y su aproximación continua (Normal). Ahora exploraremos cómo estas herramientas se integran en aplicaciones modernas de ciencia de datos: ajuste de distribuciones por máxima verosimilitud, modelado empírico de datos de conteo, y diagnóstico de sobredispersión.
  \end{retoclase}

  \vspace{0.08cm}
  \pause
  \begin{alertblock}{Preparación Recomendada}
    \begin{itemize}\itemsep=0.06cm
      \item Repasa el concepto de estimación por máxima verosimilitud (MLE) en su forma general.
      \item Reflexiona sobre cómo decidir entre Poisson y Binomial Negativa para datos de conteo empíricos.
    \end{itemize}
  \end{alertblock}
\end{frame}

\end{document}
% Archivo legado archivado: aproximación normal es subsección, no sección activa.
